\section{为什么选择视频生成：高价值模态与未成熟窗口}

上一章讲应用机会来自产品内生数据，本章看 ONE2X 为什么选择视频生成。王冠选择视频，不是因为它最容易，而是因为它价值高、模态尚未成熟、商业化空间明显，而且视频处理不只是生成单个视频，还可以形成一套围绕 DSL、context 和生产流程的系统。

\begin{figure}[H]
\centering
\includegraphics[width=0.96\textwidth]{video-generation-bet.png}
\caption{为什么选择视频生成：高价值模态、未成熟窗口、产品空间、收入验证和平台重排。自制概念图，依据 01:05:36--01:26:15 对谈内容整理。}
\end{figure}

\begin{knowledgebox}{读图：视频生成不是“做一个生成按钮”}
图里同时出现高价值、未成熟、处理空间、收入验证和平台重排。王冠看重的是视频作为信息商品的高价值，以及生成系统对内容生产、创作者工作流和平台分配权的长期影响。
\end{knowledgebox}

\subsection{视频模态为什么有商业价值}

本节先从商业价值解释视频，而不是从技术炫技解释视频。对于早期团队来说，模态选择首先要回答：用户为什么愿意付费、什么场景能最早验证价值。

视频比文本和图像更重，制作成本更高，消费带宽更大，也更贴近平台分发和商业转化。即使在生成技术不成熟时，视频处理工具和 SaaS 已经能做出可观收入。对一个早期团队来说，选择高价值模态意味着更早验证付费和创作者价值。

\subsection{Mideo 的产品判断}

接下来从模态选择进入产品形态。Mideo 的关键不是“能不能生成一段视频”，而是能否把某类内容的生产方法沉淀下来。

节目中反复强调：ONE2X 做的不是单纯 generate video，而是更广义的视频处理和生产系统。产品要理解某个品类的 production method，把它抽象成 DSL 和 recipe，再让模型生成符合目标的连续内容。王冠提到团队甚至先用自己的产品赚钱，证明它能支撑创作者生产和平台变现。

\begin{importantbox}{从生成能力到生产系统}
生成一个视频只是模型能力；稳定生产某一类视频、复用方法、迭代 recipe、形成商业闭环，才是生产系统。
\end{importantbox}

\subsection{本章小结}

视频生成的机会来自高价值模态和未成熟窗口，但真正壁垒不在“能生成”，而在能否把视频生产方法抽象为可复用系统，并在使用中产生产品内生数据。

